\section{Standards Boundary and Limitations}
\label{sec:limitations}
This section states precisely what is standards-aligned, what is prototype-specific, and what the reported evidence does and does not support. It is deliberately detailed: an implementation appendix is only useful if its deviations are enumerated.

\subsection{What follows the standards}
S-NSSAI, DNN, registration, and PDU-session procedures use the existing 3GPP mechanisms~\cite{3gpp23501}; no NAS, NGAP, RRC, or PHY message is redefined, and the S-NSSAI that the scheduler matches on is the one the core signalled for the UE's data radio bearer. The E2 path uses E2AP~v2 with the FlexRIC E2SM-RC service model (short name \texttt{ORAN-E2SM-RC}, OID \texttt{1.3.6.1.4.1.53148.1.1.2.3}) and ASN.1/PER encoding, and the added control reuses that stock encoder unmodified. The control message reproduces the E2SM-RC RRM Policy Ratio List parameter hierarchy with the standard RAN parameter identifiers. E2SM-KPM~v2.03 measurements are collected on a separate, unmodified subscription and kept as an audit record.

\subsection{Prototype-specific deviations}
The system is a research prototype, not a certified O-RAN product, and the following deviations are known and intentional.
\begin{itemize}
  \item Control Service Style~2, Action~6 (\emph{Slice-level PRB quota}) is added locally to the OAI RC RAN function. Its semantics have not undergone O-RAN conformance testing.
  \item SST and SD are transported as ASCII-text octet strings (or plain integers) rather than the 3GPP binary octet encoding, and an absent SD is defaulted to \texttt{0xffffff} by the decoder.
  \item The PLMN Identity element is required to be present in each RRM Policy Member but is never read or checked by the gNB; policy matching uses only SST and SD.
  \item The RIC Control Acknowledge is returned with a success outcome even when the gNB rejects the message. A rejection is observable only as an assertion trace and an error line in the gNB log, so a controller cannot distinguish ``applied'' from ``parsed but refused'' over E2. The measured acknowledgement rate reported in the results therefore certifies the E2 round trip, not the RAN-side acceptance.
  \item The policy service that carries intent and guidance is \emph{A1-like}: it fills the policy role this prototype needs, with versioned atomic activation, but it is not an O-RAN A1-PMS implementation~\cite{oran_a1}.
  \item The rich near-RT state is derived from FlexRIC internal service models, not from E2SM-KPM. KPM is an independent audit path and must not be presented as the source of every state feature.
  \item The core profile has slice-aware subscriber, DNN, AMF, and SMF configuration but no separately deployed NSSF. RNTI-to-UE-to-S-NSSAI attribution is reconstructed from runtime logs and configuration with an explicit validity mask; a production deployment should use a standardized identity-correlation path.
  \item One change is unrelated to the system and exists purely so the reference testbed builds: the RFSimulator option table was rewritten from designated-initializer macros to small by-value constructor helpers, because the file is compiled as C++ and the macros are rejected by the reference GCC~10 toolchain. Parameter names, defaults, and semantics are unchanged. Omitting this change makes the reproduction instructions fail.
\end{itemize}

\subsection{Scheduler limitations}
The dedicated ratio is enforced as a per-slot budget on \emph{new-data} allocation only, computed as $\lfloor n_{\mathrm{rb,avail}}\cdot r_{\mathrm{dedicated}}/100\rfloor$ over the full cell bandwidth. Several consequences follow and bound the generality of the results.
\begin{itemize}
  \item HARQ retransmissions and the small allocations for timing-advance or beam-switch control elements are served before any slice budget and are not charged against it. Slice quotas can therefore never suppress retransmissions, but retransmission PRBs do reduce the free blocks the budgeted pass can subsequently find, so the realized share can fall below the commanded share under heavy retransmission.
  \item Minimum and maximum ratios are validated at the E2 boundary but are not independent scheduler mechanisms; only the dedicated ratio is enforced, and bound projection is performed by the near-RT controller before transmission.
  \item Unused budget of one controlled slice is not lent to another controlled slice. Saturating traffic on both slices is therefore a precondition for the commanded ratio to be observable, which is why the campaign offers 150~Mbps against a cell that cannot carry it.
  \item Each UE's budgeted allocation is a single largest contiguous free block, capped by the remaining budget. Fragmentation can leave part of a slice budget unusable within a slot.
  \item A slice whose dedicated ratio is zero is skipped by its own pass and is excluded from the unmatched pass as well, so a zero dedicated ratio yields zero new-data PRBs rather than best-effort service. The operational 10\% floor prevents this in the campaign.
  \item Matching is exact on the pair (SST, SD): there is no wildcard or partial match, and a UE with no data radio bearer defaults to SST~0 / SD \texttt{0xffffff} and matches nothing. For a UE with several bearers on different slices the candidate takes the S-NSSAI of the first bearer reporting a non-zero 5QI; this is unambiguous here because every UE has exactly one PDU session and one bearer.
  \item The per-slice budget sums the available resource blocks over all beams while the allocation hook is invoked once per beam group. This is exact only for a single beam, which is the configuration used.
  \item Independently of any policy, the downlink control-channel budget limits the testbed to at most four scheduled UEs per downlink slot out of five, so no PRB policy can serve all five UEs in the same slot.
  \item The policy installer clears the table before re-validating each entry and does not restore the previous table if an entry is rejected mid-loop. This is unreachable in the deployed path because the E2 decoder applies identical bounds first and installs nothing on failure, but it is a latent fault for any other caller.
\end{itemize}
Results therefore characterize this implementation, not all possible slicing schedulers.

\subsection{Measurement limitations}
Two reported quantities are weaker than they appear. Under the static AWGN RFSimulator profile used for the pilot, the MAC block-error-rate and wideband-CQI counters are effectively unpopulated: their run means are $5.6\times10^{-45}$ and $0.0$ respectively across all three runs. They must not be cited as measured radio quality, and the BLER term of the reward is consequently inert in this pilot, so its weight is untested. Second, the reported downlink loss is an offered-versus-delivered ratio computed from the UE tunnel-interface byte counters against the nominal offered rate; it is not the sender's own datagram-loss report, and it should not be read as a control failure under an intentionally saturating offered load.

\subsection{Control-loop and learning limitations}
\begin{itemize}
  \item \emph{Guidance acted mainly as a constraint.} In the pilot, 90.66\% of guided training actions and 92.1\%/78.2\% of its two evaluation phases sat within half a PRB of an edge of the active feasible interval, against 0.34\%/0\%/0\% for the unguided arm. Because the Intent-1 guidance floor of 55\% maps to $b_A\geq 59$ and the guided evaluation mean was 59.27 PRBs, the guided policy was effectively pinned to the guidance boundary rather than exploring inside it. The guided arm's advantage in this run should therefore be attributed primarily to the calibrated constraint, not to a richer learned policy.
  \item \emph{Accepted-Actor counts are an upper bound.} The learner contains a force-promotion escape hatch that publishes a finite, legally projected candidate after six consecutive holdout rejections. With 195 rejections among 257 candidates in the guided run, the reported 62 acceptances is an upper bound on the number that passed the quality gates outright; the accepted-versus-forced split lives only in the run databases.
  \item \emph{Evaluation is not fully frozen.} Freezing suspends the Actor, replay ingestion, and the observation normalization statistics, but the guidance-fusion weight and the intent-violation fallback counters continue to update during evaluation.
  \item \emph{The two learning arms do not bootstrap over identical episodes.} A transition is marked terminal when its governing policy version changes. For the guided arm that is the A1 policy version, renewed every 10~s; for the unguided arm it is the intent version, renewed only at an intent switch. The guided arm therefore learns over markedly shorter effective episodes, which is a confound in any guided-versus-unguided comparison.
  \item \emph{One logged column is inert.} In asynchronous mode the controller-side reward scaler is never updated and is not shipped in the Actor snapshot, so the recorded running standard deviation is identically unity and the recorded scaled reward equals the raw reward. Reward scaling is performed inside the learner; the controller-side columns are an audit copy at unit scale.
\end{itemize}

\subsection{Evidence and reproducibility limitations}
The pilot is a single balanced-load seed. It cannot separate policy effects from seed-specific training and runtime variation, and no statistical claim is made from it. Beyond that, five reproducibility caveats should be recorded honestly.

First, the three run databases referenced by the artifact manifest, of 1.0--1.3~GB each, together with every trained checkpoint, were written under \texttt{/tmp} and were destroyed when the host was rebooted to repair an unrelated GPU driver fault, because that directory does not survive a reboot. Every number in this paper was re-verified against the committed comma-separated metric files and against the code that produces them, but the pilot is currently reproducible from code and protocol rather than re-derivable from raw data. Run databases must be archived outside \texttt{/tmp} before any campaign is considered complete.

Second, database freshness and the RAN/RIC restart are conditional on invoking the runner with its stack-management flag; without it, runs append to a single shared database and the mapped-UE validation query, which carries no run or time filter, counts UEs across the entire file. All campaign runs must use that flag.

Third, the result generator keys runs on the tuple (scenario, arm, seed) and does not consult the manifest mode field. A checkpoint-evaluation run validates against a reduced phase set, can legitimately reach primary quality, and---being newer---would supersede the corresponding full training run, silently removing its training phase from the tables. No claim of checkpoint-evaluation isolation is made here; until the generator is fixed, checkpoint evaluations must be written to a separate results directory.

Fourth, two thresholds present in the campaign specification are recorded but never enforced: the 99th-percentile fresh-action-interval bound and the maximum summary age. Conversely, the enforced control-latency and effect-coverage bounds are hard-coded in the runner rather than specification-driven. The set of gates actually applied is the one enumerated in the methodology section, not the full key list of the specification file.

Fifth, seeding covers the Python and PyTorch generators in both the controller and the learner process and derives the channel-trace seeds, and arm order within each seed and scenario is shuffled deterministically; but NumPy, CUDA, and cuDNN determinism are not forced, and the learner advances on wall clock. Runs are therefore statistically, not bitwise, reproducible. The cost is also material: a single run takes 85--88~minutes against a 4090~s nominal phase budget, the roughly 27\% overhead being A1 guidance regeneration across the 64 intent switches, so the full 45-run campaign requires on the order of 65~hours of continuous testbed time.

Finally, the scope of the radio experiment itself is narrow: one gNB, static AWGN, RFSimulator, local network namespaces, downlink-only control, and two slices. Commercial RU operation would additionally require timing synchronization, RF calibration, DU/RU integration, and renewed real-time validation.
